\chapter{Guía de Uso del Sistema}
\label{apendice:guia_uso}

Este apéndice proporciona instrucciones detalladas para instalar, configurar y utilizar el módulo aislado de detección de semáforos.

\section{Requisitos del Sistema}

\subsection{Hardware Recomendado}

\begin{table}[H]
\centering
\begin{tabular}{|l|l|}
\hline
\textbf{Componente} & \textbf{Especificación} \\
\hline
CPU & Intel i5/i7 o AMD equivalente (multi-core) \\
RAM & Mínimo 8 GB, recomendado 16 GB \\
GPU & NVIDIA con soporte CUDA (opcional pero recomendado) \\
Almacenamiento & 5 GB libres \\
\hline
\end{tabular}
\caption{Requisitos de hardware}
\end{table}

\subsection{Software Requerido}

\begin{itemize}
    \item Python 3.7 o superior
    \item CUDA Toolkit (opcional, si se usa GPU)
    \item Sistema operativo: Linux (Ubuntu 18.04+), macOS, o Windows 10+
\end{itemize}

\section{Instalación}

\subsection{Clonar el Repositorio}

\begin{codesnippet}
\begin{verbatim}
git clone [URL_DEL_REPOSITORIO]
cd TrafficLightDetection
\end{verbatim}
\end{codesnippet}

\subsection{Crear Entorno Virtual (Recomendado)}

\begin{codesnippet}
\begin{verbatim}
# Crear entorno virtual
python3 -m venv venv

# Activar entorno
# En Linux/macOS:
source venv/bin/activate
# En Windows:
venv\Scripts\activate
\end{verbatim}
\end{codesnippet}

\subsection{Instalar Dependencias}

\begin{codesnippet}
\begin{verbatim}
pip install -r requirements.txt
\end{verbatim}
\end{codesnippet}

Las principales dependencias incluyen:
\begin{itemize}
    \item PyTorch $\geq$ 1.8.0
    \item torchvision
    \item NumPy
    \item OpenCV (cv2)
    \item Matplotlib
\end{itemize}

\subsection{Verificar Instalación}

\begin{codesnippet}
\begin{verbatim}
# Verificar PyTorch
python -c "import torch; print(torch.__version__)"

# Verificar CUDA (si aplica)
python -c "import torch; print(torch.cuda.is_available())"
\end{verbatim}
\end{codesnippet}

\section{Uso Básico}

\subsection{API Principal}

El punto de entrada principal es la función \texttt{load\_pipeline()}:

\begin{codesnippet}
\begin{verbatim}
from tlr.pipeline import load_pipeline
import cv2

# Cargar pipeline
# Para GPU:
pipeline = load_pipeline('cuda:0')
# Para CPU:
# pipeline = load_pipeline('cpu')

# Cargar imagen
image = cv2.imread('path/to/image.jpg')

# Definir projection boxes
# Formato: [x1, y1, x2, y2] para cada ROI
projection_boxes = [
    [100, 50, 150, 120],  # box 1
    [500, 80, 550, 150],  # box 2
]

# Ejecutar detección
result = pipeline(image, projection_boxes)

# Resultado contiene:
# - result['bboxes']: coordenadas de semáforos detectados
# - result['colors']: estado de cada semáforo
# - result['scores']: confianza de cada detección
\end{verbatim}
\end{codesnippet}

\subsection{Procesamiento de Video}

\begin{codesnippet}
\begin{verbatim}
import cv2
from tlr.pipeline import load_pipeline

# Cargar pipeline
pipeline = load_pipeline('cuda:0')

# Abrir video
cap = cv2.VideoCapture('video.mp4')
frame_count = 0

while cap.isOpened():
    ret, frame = cap.read()
    if not ret:
        break

    # Procesar frame
    # (projection_boxes debería venir de HD-map o archivo)
    result = pipeline(frame, projection_boxes)

    # Visualizar resultados
    for bbox, color, score in zip(result['bboxes'],
                                   result['colors'],
                                   result['scores']):
        x1, y1, x2, y2 = bbox
        cv2.rectangle(frame, (x1, y1), (x2, y2), (0, 255, 0), 2)
        cv2.putText(frame, f'{color} ({score:.2f})',
                   (x1, y1-10), cv2.FONT_HERSHEY_SIMPLEX,
                   0.5, (0, 255, 0), 2)

    # Mostrar frame
    cv2.imshow('Traffic Light Detection', frame)
    if cv2.waitKey(1) & 0xFF == ord('q'):
        break

    frame_count += 1

cap.release()
cv2.destroyAllWindows()
\end{verbatim}
\end{codesnippet}

\section{Scripts Principales}

A continuación se describen los scripts principales utilizados para el flujo de trabajo del sistema.

\subsection{Extracción de Frames}

\begin{codesnippet}
\begin{verbatim}
python frames_extractor.py
\end{verbatim}
\end{codesnippet}

Extrae frames de un video para su posterior procesamiento. El video de entrada se selecciona directamente en el código del script, modificando la variable correspondiente al path del archivo de video.

\subsection{Selección de Projection Boxes}

\begin{codesnippet}
\begin{verbatim}
python select_projection_and_append.py
\end{verbatim}
\end{codesnippet}

Herramienta interactiva para seleccionar las projection boxes (ROIs donde se espera encontrar semáforos) manualmente usando el mouse. Las coordenadas seleccionadas se guardan en un archivo de texto para su uso posterior.

\subsection{Propagación de Projection Boxes}

\begin{codesnippet}
\begin{verbatim}
python propagate.py
\end{verbatim}
\end{codesnippet}

Propaga las projection boxes definidas en un frame a todos los frames extraídos del video. Útil cuando las projection boxes son válidas para toda la secuencia de frames.

\subsection{Ejecución del Pipeline}

\begin{codesnippet}
\begin{verbatim}
python run_pipeline.py
\end{verbatim}
\end{codesnippet}

Ejecuta el pipeline completo de detección de semáforos sobre los frames extraídos, utilizando las projection boxes definidas. Genera los resultados de detección y clasificación para cada frame.

\section{Estructura de Directorios del Proyecto}

\begin{verbatim}
TrafficLightDetection/
+-- src/
|   +-- tlr/
|       +-- pipeline.py           # Pipeline principal
|       +-- detector.py           # Detector de semaforos
|       +-- recognizer.py         # Clasificadores
|       +-- tracking.py           # Tracking temporal
|       +-- selector.py           # Asignacion hungara
|       +-- hungarian_optimizer.py
|       +-- utils.py              # Utilidades
|       +-- weights/              # Modelos pre-entrenados
|       +-- confs/                # Configuraciones
+-- inputTesis/                   # Videos de entrada
+-- docs/                         # Documentacion
|   +-- tesis/                    # Documentos de la tesis
|   +-- ...
+-- requirements.txt              # Dependencias
+-- README.md                     # Readme del proyecto
\end{verbatim}
